\chapter{Đặt vấn đề}

\section{Bối cảnh và lý do chọn đề tài}

\subsection{Bối cảnh xã hội}

Ngôn ngữ ký hiệu là phương tiện giao tiếp quan trọng đối với người khiếm
thính. Bên cạnh người khiếm thính, người thân, bạn bè, giáo viên, tình nguyện
viên và nhân viên hỗ trợ xã hội cũng có nhu cầu học các ký hiệu cơ bản để giao
tiếp thuận tiện và bao trùm hơn.

Khác với việc học từ vựng thông thường, việc học ngôn ngữ ký hiệu phụ thuộc
nhiều vào yếu tố trực quan. Người học cần quan sát hình dạng bàn tay, hướng
lòng bàn tay, vị trí tay so với cơ thể, chuyển động, nhịp độ và một phần biểu
cảm khuôn mặt. Nếu chỉ sử dụng mô tả bằng văn bản hoặc hình ảnh tĩnh, người
học có thể hiểu nghĩa của từ nhưng vẫn thực hiện sai động tác.

Nền tảng web có khả năng hỗ trợ quá trình tự học nhờ ưu điểm dễ tiếp cận trên
nhiều thiết bị và không yêu cầu cài đặt phần mềm chuyên biệt. Khi kết hợp video
minh họa, bài kiểm tra, cơ chế ôn tập và phản hồi từ camera, hệ thống có thể tạo
ra trải nghiệm học tập chủ động hơn so với việc chỉ xem video riêng lẻ.

Trong phạm vi phiên bản hiện tại, đề tài tập trung vào từ vựng American Sign
Language (ASL). Hệ thống chưa hỗ trợ Vietnamese Sign Language (VSL) hoặc dịch
hội thoại ký hiệu liên tục.

\subsection{Vấn đề từ góc nhìn tương tác người -- máy}

Từ góc nhìn Human-Computer Interaction (HCI), một hệ thống học ASL cần giải
quyết ba nhóm vấn đề chính: nhận thức, thao tác và phản hồi.

Về nhận thức, giao diện cần ưu tiên video và hình ảnh trực quan, đồng thời hạn
chế hiển thị quá nhiều thông tin cùng lúc. Nội dung nên được chia theo cấp độ
và chủ đề để người mới không phải tiếp cận một danh sách từ vựng quá lớn ngay
từ đầu.

Về thao tác, người học cần dễ dàng tìm bài học, xem lại video, bắt đầu bài kiểm
tra, tra cứu từ và bật camera để luyện tập. Các tác vụ chính cần được thể hiện
rõ ràng, tránh yêu cầu người dùng thực hiện nhiều bước không cần thiết.

Về phản hồi, hệ thống cần cho người học biết trạng thái hiện tại: video đang
phát hay tạm dừng, bài học đã hoàn thành đến đâu, camera đã sẵn sàng hay chưa
và kết quả nhận diện có độ tin cậy như thế nào. Khi xảy ra lỗi, giao diện cần
đưa ra hướng dẫn để người dùng có thể tiếp tục thao tác.

Các yêu cầu trên là cơ sở để xác định nhóm người dùng, thiết kế luồng tương tác
và xây dựng tiêu chí đánh giá trải nghiệm trong các phần tiếp theo.

\section{Đối tượng người dùng và nhu cầu}

\subsection{Người mới bắt đầu học ASL}

Đây là nhóm người dùng chính của hệ thống. Họ cần một lộ trình học dễ tiếp cận,
nội dung được chia nhỏ và video minh họa rõ ràng. Sau mỗi bài học, người dùng
cần thực hiện bài kiểm tra để đánh giá khả năng ghi nhớ và ôn lại các từ chưa
thành thạo.

\subsection{Người có nhu cầu giao tiếp cơ bản}

Nhóm này bao gồm người thân, bạn bè, giáo viên hoặc tình nguyện viên muốn học
một số ký hiệu thông dụng. Họ cần chức năng tra cứu nhanh, bài học theo chủ đề
và khả năng xem lại động tác nhiều lần.

\subsection{Quản trị viên nội dung}

Quản trị viên cần công cụ để thêm, sửa, kiểm tra và xuất bản nội dung từ vựng,
video, thông tin mô tả và câu hỏi trắc nghiệm. Bên cạnh đó, dữ liệu phân tích
giúp quản trị viên phát hiện bài học có tỷ lệ hoàn thành thấp hoặc ký hiệu
thường xuyên gặp lỗi khi nhận diện.

\section{Hạn chế của các phương pháp hiện có}

\subsection{Học qua tài liệu tĩnh và video riêng lẻ}

Tài liệu và video minh họa giúp người học quan sát động tác mẫu nhưng chưa tạo
thành một vòng học hoàn chỉnh. Người học vẫn phải tự ghi nhớ, tự đánh giá kết
quả và tự lên lịch ôn tập.

Phương pháp này tồn tại một số hạn chế:

\begin{itemize}
  \item Thiếu lộ trình học rõ ràng.
  \item Thiếu cơ chế kiểm tra sau bài học.
  \item Không theo dõi được tiến độ cá nhân.
  \item Không cung cấp phản hồi trực tiếp để hỗ trợ người học tự đánh giá động
        tác khi thực hành.
  \item Khó xác định từ nào đã thành thạo và từ nào cần ôn lại.
\end{itemize}

\subsection{Ứng dụng từ điển và nền tảng học ký hiệu}

Một số nền tảng hiện có cung cấp từ điển ký hiệu, video hoặc tài liệu theo chủ
đề. Tuy nhiên, không phải hệ thống nào cũng tích hợp đồng thời lộ trình học,
bài kiểm tra, cơ chế ôn tập ngắt quãng, bảng theo dõi tiến độ và luyện tập bằng
camera.

Bảng~\ref{tab:comparison} trình bày sự khác biệt giữa SmartSignLanguage và một
số nền tảng hỗ trợ học ngôn ngữ ký hiệu phổ biến.

\begingroup
\small
\setlength{\tabcolsep}{3pt}
\renewcommand{\arraystretch}{1.15}
\begin{longtable}{
  |>{\raggedright\arraybackslash}p{1.8cm}
  |>{\raggedright\arraybackslash}p{3.15cm}
  |>{\raggedright\arraybackslash}p{3.15cm}
  |>{\raggedright\arraybackslash}p{3.15cm}
  |>{\raggedright\arraybackslash}p{3.15cm}|}
  \caption{So sánh một số giải pháp hỗ trợ học ngôn ngữ ký hiệu}
  \label{tab:comparison}\\
  \hline
  \textbf{Tiêu chí} &
  \textbf{HandSpeak} &
  \textbf{Signing Savvy} &
  \textbf{Spread the Sign} &
  \textbf{SmartSignLanguage} \\
  \hline
  \endfirsthead

  \multicolumn{5}{c}{\tablename\ \thetable\ -- tiếp theo} \\
  \hline
  \textbf{Tiêu chí} &
  \textbf{HandSpeak} &
  \textbf{Signing Savvy} &
  \textbf{Spread the Sign} &
  \textbf{SmartSignLanguage} \\
  \hline
  \endhead

  Video minh họa &
  Có. Cung cấp từ điển ASL, video và nội dung về ngữ pháp, văn hóa Deaf. &
  Có. Cung cấp video ký hiệu và cho phép xem lại nhiều lần. &
  Có. Là từ điển trực quan hỗ trợ tra cứu ký hiệu thuộc nhiều ngôn ngữ. &
  Có. Video được tích hợp trực tiếp vào thẻ từ vựng và bài học. \\
  \hline

  Lộ trình học &
  Hạn chế. Có bài hướng dẫn được phân loại nhưng người học chủ yếu tự lựa chọn
  nội dung. &
  Hạn chế. Có danh sách từ theo chủ đề nhưng thiên về công cụ hỗ trợ tự học. &
  Không phải chức năng chính. Nền tảng tập trung vào tra cứu từ và cụm từ. &
  Có. Gồm ba cấp độ và 16 bài học được sắp xếp từ cơ bản đến nâng cao. \\
  \hline

  Bài kiểm tra &
  Hạn chế. Có một số bài luyện tập nhưng chưa phải bài kiểm tra sau từng bài
  học. &
  Có. Người dùng có thể tạo bài kiểm tra từ danh sách từ. &
  Không phải chức năng chính. &
  Có. Bài kiểm tra được thiết kế theo nội dung bài học và kết quả được lưu lại. \\
  \hline

  Ôn tập từ yếu &
  Hạn chế. Người học chủ yếu tự lựa chọn nội dung cần xem lại. &
  Hạn chế. Có thẻ ghi nhớ và danh sách từ nhưng người dùng cần chủ động lựa
  chọn nội dung ôn tập. &
  Hạn chế. Tập trung vào tra cứu hơn là quản lý quá trình ôn tập cá nhân. &
  Có. Hệ thống đưa từ yếu hoặc đến hạn vào danh sách ôn tập theo cơ chế SRS. \\
  \hline

  Theo dõi tiến độ &
  Hạn chế. Không tập trung vào bảng theo dõi tiến độ cá nhân. &
  Hạn chế. Có công cụ hỗ trợ tự học nhưng không tập trung vào một lộ trình
  thống nhất. &
  Hạn chế. Không tập trung vào quản lý tiến độ học tập. &
  Có. Bảng theo dõi tổng hợp tiến độ bài học, kết quả bài kiểm tra và hoạt động
  luyện tập. \\
  \hline

  Luyện tập bằng camera &
  Không. Người học xem video và tự đối chiếu động tác. &
  Không. Người học luyện tập thông qua video, thẻ ghi nhớ và bài kiểm tra. &
  Không. Nền tảng cung cấp video để tra cứu và đối chiếu. &
  Có. AI nhận diện ký hiệu qua camera hoặc tệp tải lên và đưa ra phản hồi sơ bộ
  trong phạm vi tập nhãn mà mô hình đã được huấn luyện. \\
  \hline
\end{longtable}
\endgroup

Qua bảng so sánh, HandSpeak có thế mạnh về tài nguyên ASL, nội dung ngữ pháp và
văn hóa Deaf. Signing Savvy hỗ trợ việc tự học thông qua danh sách từ, thẻ ghi
nhớ và bài kiểm tra. Spread the Sign phù hợp với nhu cầu tra cứu ký hiệu thuộc
nhiều ngôn ngữ khác nhau.

SmartSignLanguage không nhằm thay thế hoàn toàn các nền tảng trên. Điểm khác
biệt của đề tài là tích hợp các chức năng học tập vào một luồng sử dụng thống
nhất: chọn bài học, xem video, thực hiện bài kiểm tra, ôn tập từ yếu, theo dõi
tiến độ và luyện tập bổ sung bằng camera.

\section{Khoảng trống cần giải quyết}

Từ các hạn chế trên, đề tài xác định bốn khoảng trống chính:

\begin{itemize}
  \item Thiếu một môi trường học ASL tích hợp bài học, video, bài kiểm tra và
        bảng theo dõi tiến độ trong một luồng sử dụng thống nhất.
  \item Thiếu cơ chế ôn tập ngắt quãng (Spaced Repetition System -- SRS) để tự
        động xác định các từ yếu hoặc đến hạn ôn tập.
  \item Thiếu cơ chế phản hồi sơ bộ khi người học thực hành bằng camera.
  \item Thiếu hệ thống quản trị nội dung linh hoạt để cập nhật, kiểm tra dữ liệu
        động và theo dõi dữ liệu phân tích liên quan đến bài học, câu hỏi trắc
        nghiệm và kết quả nhận diện.
\end{itemize}

SmartSignLanguage được xây dựng nhằm giải quyết các khoảng trống này trong
phạm vi một hệ thống web hỗ trợ học ASL ở mức từ vựng và cụm từ cơ bản.

\section{Mục tiêu của đề tài}

\subsection{Mục tiêu tổng quát}

Xây dựng một ứng dụng web hỗ trợ học và luyện tập từ vựng ASL theo hướng trực
quan, dễ sử dụng và có khả năng theo dõi tiến độ cá nhân.

\subsection{Mục tiêu chức năng}

Hệ thống hướng đến các chức năng chính:

\begin{itemize}
  \item Cung cấp lộ trình học gồm ba cấp độ: Beginner, Intermediate và Advanced.
  \item Tổ chức 16 bài học theo các chủ đề như Greetings, Family, Colors, Food
        \& Drink, School, Travel và Deaf Community.
  \item Hiển thị thẻ từ vựng kèm video, mô tả và hướng dẫn luyện tập.
  \item Cung cấp bài kiểm tra sau bài học.
  \item Hỗ trợ ôn tập các từ yếu hoặc đến hạn theo cơ chế SRS.
  \item Cho phép tra cứu ký hiệu.
  \item Nhận diện ký hiệu trong phạm vi mô hình đã huấn luyện: camera hỗ trợ
        Words, Alphabet và Numbers; tệp tải lên hỗ trợ Alphabet và Numbers.
  \item Chuyển đổi văn bản đơn giản thành chuỗi hoạt ảnh điểm mốc cơ thể.
  \item Cung cấp bảng theo dõi tiến độ cho người học.
  \item Cung cấp công cụ quản trị để thêm, sửa, kiểm tra và xuất bản dữ liệu từ
        vựng, video và câu hỏi trắc nghiệm.
  \item Cung cấp dữ liệu phân tích để hỗ trợ quản trị viên phát hiện bài học có
        tỷ lệ hoàn thành thấp hoặc ký hiệu có tỷ lệ nhận diện thất bại cao.
\end{itemize}

\subsection{Mục tiêu HCI}

Bên cạnh việc hoàn thiện chức năng, đề tài tập trung vào các mục tiêu HCI:

\begin{itemize}
  \item Giảm số bước cần thiết để người dùng bắt đầu một bài học.
  \item Giúp người dùng dễ dàng tìm kiếm và xem lại nội dung.
  \item Hiển thị trạng thái rõ ràng khi tải dữ liệu, sử dụng camera hoặc nhận
        kết quả AI.
  \item Hỗ trợ người dùng phục hồi khi thao tác sai hoặc khi camera không hoạt
        động.
  \item Hạn chế quá tải nhận thức bằng cách chia nội dung thành cấp độ và chủ
        đề.
  \item Đảm bảo các chức năng học cơ bản vẫn sử dụng được khi người dùng không
        cấp quyền camera.
\end{itemize}

\section{Đề xuất giải pháp SmartSignLanguage}

SmartSignLanguage là ứng dụng web hỗ trợ học, tra cứu và luyện tập ASL. Thay vì
chỉ cung cấp video riêng lẻ, hệ thống xây dựng một vòng học gồm các bước:

\begin{enumerate}
  \item Người dùng chọn bài học theo cấp độ và chủ đề.
  \item Người dùng xem video và thông tin của từng thẻ từ vựng.
  \item Người dùng hoàn thành bài kiểm tra để đánh giá khả năng ghi nhớ.
  \item Hệ thống lưu tiến độ và đưa các từ yếu vào danh sách ôn tập.
  \item Người dùng có thể thực hành bổ sung bằng camera.
  \item Bảng theo dõi hiển thị tiến độ và kết quả học tập.
\end{enumerate}

Chức năng nhận diện sử dụng MediaPipe để trích xuất điểm mốc cơ thể (landmark)
của bàn tay và tư thế từ các khung hình. Các tọa độ landmark được chuẩn hóa và
tổng hợp thành chuỗi dữ liệu đầu vào. Sau đó, mô hình học sâu phân loại ký hiệu
trong tập nhãn đã được huấn luyện. Kết quả hiển thị gồm nhãn ký hiệu dự đoán, độ
tin cậy và phản hồi sơ bộ hỗ trợ luyện tập.

Phản hồi từ AI không thay thế giáo viên. Hệ thống chưa phân tích chính xác từng
lỗi nhỏ về hình dạng bàn tay, hướng tay, chuyển động hoặc biểu cảm khuôn mặt.

Chức năng Text-to-Sign chuyển đổi văn bản đơn giản thành chuỗi hoạt ảnh
landmark. Hệ thống ưu tiên tìm từ hoặc cụm từ có sẵn trong dữ liệu. Nếu không
tìm thấy, hệ thống có thể biểu diễn từ bằng cách đánh vần ngón tay
(fingerspelling) khi dữ liệu chữ cái khả dụng. Đây chưa phải cơ chế dịch hoàn
chỉnh từ ngôn ngữ tự nhiên sang ngữ pháp ASL.

\section{Phạm vi và giới hạn}

\subsection{Phạm vi}

Phiên bản hiện tại tập trung vào:

\begin{itemize}
  \item Học ASL ở mức từ vựng và một số cụm từ cơ bản.
  \item Ba cấp độ học và 16 bài học theo chủ đề.
  \item Bài kiểm tra, ôn tập và theo dõi tiến độ.
  \item Nhận diện Words, Alphabet và Numbers trong phạm vi mô hình.
  \item Chuyển đổi văn bản đơn giản thành hoạt ảnh landmark.
  \item Quản trị nội dung và theo dõi dữ liệu phân tích.
\end{itemize}

\subsection{Giới hạn}

Hệ thống hiện còn một số giới hạn:

\begin{itemize}
  \item Chưa hỗ trợ VSL.
  \item Chưa hỗ trợ nhận diện hoặc dịch câu ký hiệu liên tục.
  \item Tập từ vựng AI còn giới hạn bởi tập dữ liệu và mô hình đã huấn luyện.
  \item Kết quả nhận diện phụ thuộc vào ánh sáng, vị trí camera và tốc độ chuyển
        động.
  \item Text-to-Sign chưa xử lý đầy đủ ngữ pháp ASL.
  \item Hoạt ảnh landmark chưa tự nhiên như video người thật hoặc avatar 3D có
        biểu cảm.
  \item Recognition là chức năng luyện tập bổ sung, chưa phải bước bắt buộc
        trong mỗi bài học.
\end{itemize}

Các giới hạn này cần được trình bày rõ để tránh hiểu nhầm về năng lực của hệ
thống và làm cơ sở cho hướng phát triển sau này.

\section{Định hướng đánh giá trải nghiệm người dùng}

Để đánh giá hệ thống từ góc nhìn HCI, báo cáo dự kiến sử dụng các tiêu chí:

\begin{itemize}
  \item Tỷ lệ hoàn thành các tác vụ cơ bản.
  \item Thời gian trung bình để tìm và bắt đầu một bài học.
  \item Thời gian trung bình để tra cứu một ký hiệu.
  \item Số lỗi thao tác khi bật camera và bắt đầu recognition.
  \item Mức độ dễ sử dụng thông qua bảng hỏi Likert hoặc System Usability Scale
        (SUS).
  \item Mức độ hài lòng đối với video, bài kiểm tra, danh sách ôn tập và phản
        hồi recognition.
  \item Nhận xét định tính của người dùng sau khi trải nghiệm hệ thống.
\end{itemize}

Các chỉ số HCI cần được đánh giá riêng với chỉ số kỹ thuật của AI như độ chính
xác, độ tin cậy hoặc ma trận nhầm lẫn.

\section{Cấu trúc báo cáo}

Nội dung báo cáo được tổ chức như sau:

\begin{itemize}
  \item \textbf{Chương I:} Trình bày bối cảnh, đối tượng người dùng, hạn chế của
        các giải pháp hiện có, mục tiêu, phạm vi và giải pháp đề xuất.
  \item \textbf{Chương II:} Trình bày kiến thức nền tảng về ASL, HCI, thị giác
        máy tính, MediaPipe, học sâu và cơ chế ôn tập ngắt quãng.
  \item \textbf{Chương III:} Thu thập, phân tích và đặc tả yêu cầu hệ thống.
  \item \textbf{Chương IV:} Trình bày thiết kế giao diện, kiến trúc hệ thống, cơ
        sở dữ liệu, API và luồng xử lý.
  \item \textbf{Chương V:} Trình bày quá trình cài đặt, kết quả triển khai, kiểm
        thử và đánh giá trải nghiệm người dùng.
  \item \textbf{Chương VI:} Tổng kết kết quả đạt được, giới hạn và hướng phát
        triển.
\end{itemize}
